\chapter*{Introduction Générale}
\addcontentsline{toc}{chapter}{Introduction Générale}
\markboth{Introduction Générale}{Introduction Générale}

\begin{infobox}[Contexte général]
Les systèmes de santé font face à une pression croissante~: vieillissement de la
population, multiplication des maladies chroniques et explosion des coûts
hospitaliers. Selon l'OCDE, les hospitalisations évitables représentent jusqu'à
\textbf{15\,\% des dépenses hospitalières} dans les pays développés, soit plusieurs
milliards de dollars chaque année.
\end{infobox}

\vspace{1em}

\noindent L'émergence du \textbf{Machine Learning Operations (MLOps)} offre une
réponse concrète à ce défi~: industrialiser les modèles d'intelligence artificielle
pour qu'ils produisent une valeur réelle, continue et fiable en production, et non
uniquement dans des notebooks de recherche. Ce projet se situe précisément à
l'intersection de la santé numérique et du MLOps.

\section*{Contexte et problématique}

Le dataset CMS~DE-SynPUF (Centers for Medicare \& Medicaid Services) constitue
l'une des bases de données médicales synthétiques les plus complètes et réalistes
disponibles publiquement. Il couvre 2,3~millions de bénéficiaires sur la période
2008--2010, avec plus de 60~Go de données incluant les claims d'hospitalisation,
de consultations ambulatoires, de prescriptions et de données démographiques.

\noindent La problématique centrale de ce projet est la suivante~:

\begin{criticalbox}[Question de recherche]
\textit{Comment construire un système de prédiction du risque d'hospitalisation
suffisamment fiable pour guider des interventions préventives, tout en garantissant
que ce système reste opérationnel, équitable et traçable au fil du temps~?}
\end{criticalbox}

\noindent Cette question soulève plusieurs défis~:
\begin{itemize}
  \item le \textbf{déséquilibre de classes} (seuls 8.1\,\% des patients sont
        hospitalisés dans l'année) ;
  \item la \textbf{fuite de données temporelle} (\textit{data leakage}) si la
        variable cible est mal construite ;
  \item la \textbf{dérive des données} dans le temps, qui dégrade les performances
        sans qu'on s'en rende compte ;
  \item les \textbf{biais algorithmiques} qui pénalisent certains groupes raciaux
        et socio-économiques.
\end{itemize}

\section*{Objectifs du projet}

Ce projet vise à concevoir, entraîner, déployer et surveiller une plateforme MLOps
complète pour la détection précoce du risque d'hospitalisation. Les objectifs
spécifiques sont~:

\begin{enumerate}
  \item Réaliser une \textbf{analyse exploratoire} rigoureuse et une ingénierie
        des données anti-leakage sur le dataset CMS~DE-SynPUF.
  \item Entraîner et optimiser \textbf{trois modèles} (Régression Logistique,
        XGBoost, LightGBM) avec un split temporel strict et une optimisation
        d'hyperparamètres par Optuna.
  \item \textbf{Déployer} le meilleur modèle via une API FastAPI conteneurisée
        avec Docker.
  \item Mettre en place un \textbf{pipeline de monitoring} automatique (Evidently,
        Grafana) avec détection de dérive et rollback.
  \item Évaluer l'\textbf{équité algorithmique} par groupe racial et calculer le
        ROI financier de la solution.
\end{enumerate}

\section*{Plan du rapport}

Ce rapport est organisé en quatre chapitres~:

\begin{itemize}
  \item \textbf{Chapitre~1 — Comprendre le Problème avant de le Résoudre~:}
        définition du MLOps, contexte médical, description du dataset, stack
        technologique et architecture globale du système.

  \item \textbf{Chapitre~2 — De la Donnée Brute à la Donnée Prête~:}
        analyse exploratoire complète, nettoyage des données, construction
        anti-leakage de la variable cible, fenêtres temporelles et ingénierie
        des 29 features finales.

  \item \textbf{Chapitre~3 — Apprendre à Prédire : Modèles et Expérimentation~:}
        traitement du déséquilibre, entraînement et comparaison des modèles,
        optimisation du seuil de décision, interprétabilité SHAP, traçabilité
        MLflow et analyse ROI.

  \item \textbf{Chapitre~4 — Du Modèle au Service : Mise en Production~:}
        API FastAPI, conteneurisation Docker, CI/CD GitHub~Actions, monitoring
        de la dérive, rollback automatique, analyse de fairness et bilan des
        tests réalisés.
\end{itemize}
